\section{Incrementos presentados del proyecto}

El desarrollo se realiza mediante entregas periódicas de software utilizable. A continuación se reporta el primer incremento concluido:

\subsection{Incremento 1: Autenticación RBAC y Gestión de Capacitaciones}

El objetivo de este primer ciclo fue montar la base técnica del sistema en Laravel 13, conectar la base de datos MySQL 8.0 y permitir que los administradores gestionen cursos de forma segura.

\subsubsection{Funcionalidad desarrollada}
\begin{itemize}[nosep]
    \item \textbf{Estructura base:} Proyecto Laravel 13 con tablas y migraciones vinculadas en MySQL 8.0.
    \item \textbf{Seguridad y roles (RBAC):} Inicio de sesión con permisos diferenciados para Administrador, Docente y Participante mediante políticas (\texttt{CoursePolicy}).
    \item \textbf{Gestión de cursos:} Formulario para crear y editar cursos con código único, título, fechas de inicio/fin, horas académicas, límite de vacantes y docente asignado.
    \item \textbf{Catálogo de cursos:} Interfaz con buscador reactivo y filtros por estado (Abierto, En curso, Concluido, Cancelado).
    \item \textbf{Integración con API RENIEC / PeruDevs:} Servicio desacoplado \texttt{DniLookupService} y controlador REST que consulta en tiempo real nombres y apellidos ante PeruDevs, con caché local de 30 días para evitar latencia y consumo reiterativo.
    \item \textbf{Frontend SPA Reactivo (Inertia.js v3 + Vue 3):} Vistas dinámicas con Tailwind CSS, navegación fluida sin recargas y componente de consulta DNI en tiempo real.
\end{itemize}

\subsubsection{Arquitectura unificada del servicio de consulta DNI}
Para garantizar que la verificación de ciudadanos sea reutilizable en todos los módulos (registro, matrícula y emisión de certificados) sin duplicar código, se implementó el patrón de inversión de dependencias:
\begin{itemize}[nosep]
    \item \textbf{Objeto de Transferencia (\texttt{DniPerson}):} DTO inmutable con tipado estricto que encapsula los datos retornados por la API oficial (DNI, nombres, apellidos paterno/materno, género y código verificador).
    \item \textbf{Contrato (\texttt{DniLookupService}):} Interfaz en \texttt{app/Contracts/} que desacopla la lógica de negocio del proveedor externo.
    \item \textbf{Implementación con Resiliencia (\texttt{PeruDevsDniService}):} Servicio en \texttt{app/Services/Dni/} que pre-valida el formato de 8 dígitos, aplica un tiempo de espera controlado (15 segundos) y almacena en la caché de Laravel los resultados válidos por 30 días.
    \item \textbf{Controlador REST y Componente Vue:} Endpoint \texttt{GET /api/dni/\{dni\}} consumido asíncronamente por el componente \texttt{DniSearchCard.vue} integrado en el catálogo de capacitaciones y el panel de control.
\end{itemize}

\subsubsection{Evidencia de funcionamiento}
La plataforma ya funciona en el entorno local, permitiendo iniciar sesión y crear capacitaciones.

\paragraph{Verificación de correo electrónico.}
Como parte del flujo de autenticación, el sistema envía al usuario un correo
con un enlace temporal para confirmar su dirección de correo electrónico.
Esta evidencia demuestra que la verificación de cuentas se encuentra integrada
en el primer incremento.

\begin{figure}[H]
    \centering
    \includegraphics[width=0.9\textwidth]{captura-verificacion-correo.png}
    \caption{Correo de verificación enviado por SIGC-CUSCO.}
    \label{fig:correo-verificacion}
\end{figure}

En la \cref{fig:correo-verificacion} se observa el botón \enquote{Verify Email
Address} y el enlace alternativo que puede copiarse en el navegador. El enlace
se genera de forma temporal para evitar que la confirmación permanezca válida
indefinidamente.

\begin{cajaevidencia}{Capturas de pantalla del sistema}
    \textbf{¿Qué debe ingresar aquí?:}
    Capturas de pantalla de la aplicación web: formulario de login, formulario para crear una capacitación y tabla de cursos guardados en la base de datos.
    \vspace{0.15cm}
    \textbf{Ejemplo sugerido:}
    Pegar una o dos imágenes mostrando la pantalla de login y el listado de capacitaciones activas en \texttt{http://localhost:8000/cursos}.
\end{cajaevidencia}

\subsubsection{Pruebas o validaciones realizadas}
Se implementó una suite integral de pruebas automáticas con Pest PHP y PHPUnit sobre una base de datos MySQL dedicada (\texttt{SIGC\_testing}), garantizando el aislamiento del entorno de producción y desarrollo:

\begin{table}[H]
\centering
\small
\begin{tabularx}{\textwidth}{P{1.8cm} P{4.8cm} P{4.5cm} Y}
\toprule
\textbf{Código} & \textbf{Prueba realizada} & \textbf{Resultado esperado} & \textbf{Estado} \\
\midrule
\textbf{CP-01} & Alumno intenta entrar a la pantalla de crear cursos. & Acceso denegado (Error HTTP 403 vía \texttt{CoursePolicy}). & \textbf{Aprobado} \\
\addlinespace
\textbf{CP-02} & Se intenta guardar un curso con vacantes en 0 o negativo. & Formulario bloquea y muestra mensaje de alerta en español. & \textbf{Aprobado} \\
\addlinespace
\textbf{CP-03} & Se ingresa una fecha de fin menor a la fecha de inicio. & Error de validación por fechas incoherentes (\texttt{after\_or\_equal}). & \textbf{Aprobado} \\
\addlinespace
\textbf{CP-04} & Se completa el formulario con datos válidos y docente. & Curso persistido en base de datos MySQL y visible en catálogo. & \textbf{Aprobado} \\
\addlinespace
\textbf{CP-05} & Consulta de DNI mediante servicio externo PeruDevs. & Retorna objeto estructurado \texttt{DniPerson} y almacena en caché 30 días. & \textbf{Aprobado} \\
\bottomrule
\end{tabularx}
\caption{Pruebas automáticas ejecutadas en el Incremento 1.}
\label{tab:pruebas_inc1}
\end{table}

\paragraph{Resultado de ejecución de la suite de pruebas.}
A continuación se reporta la ejecución de la batería completa de pruebas automatizadas:

\begin{quote}
    \texttt{\$ php artisan test} \\
    \texttt{Tests\textbackslash Feature\textbackslash Auth\textbackslash RegistrationTest ................. PASS} \\
    \texttt{Tests\textbackslash Feature\textbackslash CourseTest ................................... PASS} \\
    \texttt{  [OK] no autenticado o no admin no puede acceder a crear cursos} \\
    \texttt{  [OK] capacidad debe ser un entero positivo} \\
    \texttt{  [OK] fecha de fin debe ser posterior o igual a fecha de inicio} \\
    \texttt{  [OK] admin puede registrar curso valido en mysql} \\
    \texttt{  [OK] endpoint dni retorna 422 para dni invalido} \\
    \texttt{Tests\textbackslash Feature\textbackslash Services\textbackslash DniLookupServiceTest ............ PASS} \\
    \texttt{  [OK] consulta exitosa retorna DniPerson estructurado} \\
    \texttt{  [OK] respuesta exitosa se almacena en cache 30 dias} \\
    \texttt{\textbf{Tests: 50 passed (166 assertions)}} \\
    \texttt{Duration: 2.50s}
\end{quote}

\subsubsection{Cambios realizados respecto al avance anterior}
\begin{itemize}[nosep]
    \item \textbf{Uso formal de Jira:} Se abandonaron notas externas para controlar cada tarea con tickets de Jira (\texttt{SIGC-XX}).
    \item \textbf{Tablas ordenadas en MySQL:} Se separaron las tablas de usuarios, roles y cursos para no duplicar datos y asegurar relaciones limpias.
    \item \textbf{Validación estricta en Laravel:} Se configuraron reglas automáticas para impedir que se ingresen cupos negativos o fechas erróneas.
\end{itemize}

\subsubsection{Métricas y seguimiento del proceso}
\begin{itemize}[nosep]
    \item \textbf{Tiempo por tarea (\textit{Cycle Time}):} En promedio, cada historia tomó 2 días desde que se empezó a programar hasta que se aprobó.
    \item \textbf{Control de tareas en curso:} Se respetó el tope de 2 tareas simultáneas, evitando amontonar trabajo sin terminar.
    \item \textbf{Cumplimiento:} Se terminaron al 100\% las dos historias comprometidas (\texttt{SIGC-1} y \texttt{SIGC-2}).
\end{itemize}

\begin{cajaevidencia}{Captura del tablero Jira}
    \textbf{¿Qué debe ingresar aquí?:}
    Captura de pantalla de Jira mostrando las tarjetas en la columna \textbf{Done} (Hecho).
    \vspace{0.15cm}
    \textbf{Ejemplo sugerido:}
    Pegar imagen del tablero mostrando los tickets \texttt{SIGC-1} (Roles) y \texttt{SIGC-2} (Cursos) en la columna \textbf{Done}.
\end{cajaevidencia}

\subsubsection{Acuerdos de mejora para el siguiente incremento}
\begin{enumerate}[nosep]
    \item \textbf{Revisión rápida de código:} Atender los Pull Requests en menos de 24 horas para no frenar a los compañeros.
    \item \textbf{Escribir pruebas antes del código:} Aplicar TDD para la matrícula de alumnos y control de aforo antes de diseñar la pantalla.
    \item \textbf{Alistar la lectura de QR:} Probar la librería de escaneo con cámara antes de empezar el segundo incremento.
\end{enumerate}
